% =====================================================================
% Section 8 — Conclusion and Future Work
% =====================================================================
\section{Conclusion and Future Work}\label{sec:conclusion}

We presented JxMVC, a full-stack MVC framework for Jakarta~EE that
deliberately occupies an under-served point in the design space:
\SI{253}{\kilo\byte} of zero-dependency code, with security applied by
default in the request pipeline rather than opted into per route. We
described its fifteen-stage pipeline and its from-scratch subsystems,
and evaluated it against four mainstream frameworks and a GraalVM
native build using a reproducible, isolated, single-command harness.

The evidence supports a measured claim. JxMVC produces the smallest
JVM-mode deployable image of the group, starts faster than Spring Boot
and Micronaut, and sustains competitive throughput with
indistinguishable tail latency, serving \num{92759575} requests with
zero errors and zero non-2xx responses. Its clear trade-off is higher
resident memory, inherited from the embedded servlet container. In
short, a zero-dependency stack built directly on the JDK and Jakarta
Servlet API pays no meaningful latency penalty and only a modest
throughput gap, while gaining a small, auditable, secure-by-default
core.

\paragraph{Future work.} The most impactful direction is reducing
resident memory, the design's one clear weakness. Caching request-binding
metadata would further shrink reflection on the hot path and, together
with reducing dynamic reflection, would make JxMVC a viable
GraalVM native-image target---the reference row in our evaluation
quantifies the order-of-magnitude footprint gains that would follow.
Additional directions include more OIDC providers beyond the current
client, HTTP/2 and server-sent events, and broadening the empirical
evaluation to persistence- and template-heavy workloads that stress
paths beyond baseline request handling.

\section{Software and data availability}\label{sec:availability}
JxMVC and the complete, reproducible benchmark harness are openly
available at \url{https://github.com/Andre031222/jxmvc} under the MIT
license. The repository includes the framework core, the demo
application, the one-command Docker benchmark harness with its per-run
raw data and environment description (under \texttt{benchmarks/}), and
the \LaTeX{} sources of this article (under \texttt{paper/}). The core
is verified by a \NTESTS-check test suite written without any external
testing framework. The evaluated release, v3.4.0, is additionally
archived with a citable DOI on Zenodo.\footnote{DOI to be inserted upon
archival of the v3.4.0 release.}
